% applications/railway_vs_aircraft_dynamics.tex

\chapter{Railway Dynamics versus Aircraft Dynamics}

\label{chap:railway-vs-aircraft}

This chapter provides a comparative application perspective that
clarifies railway-specific consequences of constrained motion.

% ============================================================
\section{Comparative Context}
% ============================================================

Aircraft and railway domains share dynamic principles such as inertial
response and comfort relevance. The key difference for applications is
kinematic freedom: aircraft trajectories are actively steered in free
space, whereas railway vehicles are constrained by infrastructure
geometry.

% ============================================================
\section{Railway Consequence}
% ============================================================

Because railway motion is infrastructure-constrained, alignment,
cant, and runtime interpretation become primary control levers for
operational behaviour.

In application terms, the engineering question is often how existing
geometry can support admissible operational envelopes through runtime
and realization-aware interpretation.

% ============================================================
\section{AIM Interpretation}
% ============================================================

Within AIM, this comparison motivates use of canonical pose3,
vehicle-space interpretation, and operational runtime evaluation for
railway decision support.

The comparison is explanatory only; it does not introduce a parallel
architectural model.

% ============================================================
\section{Connection to AIM}
% ============================================================

\begin{summary}
The railway-versus-aircraft comparison clarifies why railway
applications require strong coupling between infrastructure semantics,
runtime services, and operational interpretation.

It reinforces application consequences of canonical AIM concepts without
redefining them.
\end{summary}
